\section{Research Gaps}
\label{sec:RG}

Despite its prominence in policy, strategy, and academic discourse, resilience remains inconsistently operationalized and unevenly assessed in organizational settings. This challenge does not stem from a lack of theoretical development. As shown above, resilience has been extensively conceptualized and modeled across disciplines, including engineering, ecology, psychology, and complex systems. However, these contributions rest on distinct ontological assumptions and methodological traditions that rarely converge into a unified organizational framework. As resilience diffuses into managerial and regulatory domains, this plurality becomes consequential: organizations are expected to enact resilience without clear guidance on which assumptions, metrics, or governance logics should prevail.\\

As shown above, resilience has been used to denote recovery, robustness, adaptation, learning, or transformation, often simultaneously and within the same organizational setting. These meanings rest on incompatible assumptions about system behavior, temporality, and control: recovery- and robustness-oriented views presuppose a stable baseline to which systems should return, whereas adaptive and transformative views assume non-stationarity, continual change, and the possibility that post-disruption states differ fundamentally from prior ones. When imported into organizations, this conceptual plurality can produce misalignment between investment, evaluation, and governance logics: organizations may invest in robustness while being evaluated on adaptability, or implement compliance controls while being rhetorically charged with transformation.\\

Our review identifies three interlinked gaps that emerge not from theoretical absence, but from insufficient integration and translation into organizational research and practice: conceptual plurality without operational convergence, methodological fragmentation across domains, and limited empirical validation in organizational contexts. {\it Conceptually}, the absence of an integrated operational core within management and information systems research prevents cumulative theory-building. {\it Methodologically}, existing metrics and formal models are rarely embedded in organization-level governance and evaluation systems, and therefore struggle to capture resilience as a dynamic, socio-technical capability. {\it Empirically}, the lack of standardized and longitudinal indicators limits comparability and learning across cases. Among these challenges, the development of operationalizable and transferable measures remains foundational \autocite{woods2015four}, a concern repeatedly documented across domains ranging from infrastructure to healthcare \autocite{heinimann2017infrastructure, kruk2017building}. Details of the systematic review process, covering search, screening, and coding across 340 publications, are provided in Appendix~\ref{appendix:B}, while illustrative field evidence of how these tensions manifest in organizational practice is synthesized in Appendix~\ref{appendix:C}.\\
\\
First, {\bf conceptual plurality without operational convergence:} Resilience has been conceptualized across domains through lenses of mechanical stability, ecological transformation, cognitive adaptability, and structural flexibility (see Appendix \ref{appendix:A}). While each view offers valuable and often rigorously formalized insights, their assumptions are not equivalent. Engineering resilience emphasizes restoration of predefined functions; ecological resilience foregrounds regime shifts and adaptive reorganization; psychological resilience centers on coping and growth under stress; socio-technical perspectives highlight emergence and distributed coordination. Although these perspectives can be interpreted as complementary at an abstract level, their implications for organizational design, measurement, and accountability differ substantially.\\

In practice, this plurality can generate definitional drift. Terms such as resistance, robustness, recovery, agility, and adaptation are often used interchangeably, despite referring to distinct mechanisms and temporal logics. Consequently, resilience may be invoked symbolically, through audits, strategy documents, or compliance templates, rather than embedded as a coherent adaptive capability \autocite{duchek2020organizational}. Without clearer articulation of how these perspectives translate into measurable organizational mechanisms, resilience risks remaining conceptually rich yet operationally diffuse.\\
\\
Second, {\bf methodological fragmentation and limited transfer:} This gap is visible in the cross-disciplinary classification of 89 resilience-related publications (see Table \ref{tab:ResearchGaps}). Simulation models, network analyses, and quantification frameworks are well established in engineering, complex systems, and climate research \autocite{francis2014metric}. However, their systematic integration into domains central to organizational and managerial resilience -- such as cybersecurity, risk management, and strategic governance -- remains limited. Table \ref{tab:ResearchGaps} shows that \textit{Markers \& Metrics} and \textit{Technology} are comparatively underrepresented in these fields, and that simulation-based approaches are rarely incorporated into organization-level governance frameworks \autocite{linkov2013resilience, kuypers2018designing}. The result is a continued reliance on maturity models, expert judgement, or audit documentation—approaches that may support compliance but rarely capture adaptive performance under real-world conditions.\\

Even within cybersecurity, formal models of attack–defense dynamics and adaptive response strategies exist \autocite{aydin2018integration}, but are seldom integrated into broader organizational resilience architectures. Other transferable approaches, such as ecological network analysis that quantifies interdependencies \autocite{luke2012systems}, transportation frameworks combining metrics and recovery modeling \autocite{jacobs2017applying}, or agent-based methods for infrastructure risk assessment \autocite{dawson2011agentbased}, remain under-exploited in digital resilience planning. The challenge is therefore not methodological scarcity, but limited cross-domain integration: technical innovations in measurement are well developed but insufficiently translated into organization-ready governance tools \autocite{heinimann2017infrastructure, aerts2018integrating}.\\
\\
\begin{center}
\newgeometry{top=.15cm, bottom=2cm, left=2.3cm, right=2.3cm}
\singlespacing
\begin{landscape}
\tiny
\begin{longtblr}[
  caption = {\textbf{Classification of Resilience-related Publications.} This table provides the classification of 89 relevant, impactful and supporting publications associated with resilience according to their corresponding topic/-s and used methodology/-ies. This table allows to extract the key research gaps, where no publication was found (red cells) or only very few (orange cell). Topic-methodology intersections in grey show other research that are less relevant for this research.},
  label = {tab:ResearchGaps}
]{
  width=\linewidth,
  colspec={|Q[1.5cm] Q[1.7cm]|Q[2.6cm] Q[2.6cm] Q[2.6cm] Q[2.6cm] Q[2.6cm] Q[2.6cm] Q[2.6cm]|},
  rowhead = 2,
  hlines, vlines
}
\SetCell[r=2, c = 2]{l} \vspace{2em}\par \textit{Topics}& & \SetCell[r = 1, c = 7]{halign = l}{\textit{Methodologies}} \\
& & \textbf{\makecell{Simulation \\ \& \\ Modeling}} & \textbf{\makecell{Network \\ Analysis}} & \textbf{\makecell{Markers \\ \& \\ Metrics}} & \textbf{\makecell{Technology}} & \textbf{\makecell{Case Studies}} & \textbf{\makecell{Longitudinal \\ Studies}} & \textbf{\makecell{Surveys}} \\
\hline
% === Example block start ===
\SetCell[r=1, c=2]{l}{\textbf{\makecell{Complex Adaptive Systems}}} & &  \SetCell[r=1, c = 1]{l} {\cite{holling1973resilience} \\ \cite{gao2016universal} \\ \cite{pumpuni2017resilience} \\ \cite{krakovska2023resilience}} & \SetCell[r=1, c = 1]{l}{\cite{albert2000error} \\ \cite{bodin2009role} \\ \cite{gao2016universal} \\ \cite{pumpuni2017resilience}} & \SetCell[r=1, c = 1]{l} {\cite{anderies2004framework} \\ \cite{sundstrom2014transdisciplinary} \\ \cite{pumpuni2017resilience}} & \SetCell[r=1, c = 1]{bg = gray!50} & \SetCell[r=1, c = 1]{l} {\cite{ostrom1990governing} \\ \cite{lanzara1999transient} \\ \cite{folke2004regime} \\ \cite{weick2011managing}} & \SetCell[r=1, c = 1]{bg = gray!50} & \SetCell[r=1, c = 1]{l}{\cite{ostrom1990governing} \\ \cite{levin1998ecosystems} \\ \cite{gunderson2002panarchy} \\ \cite{walker2004resilience} \\ \cite{heinimann2017infrastructure}}\\
\hline
\SetCell[r=1,c=2]{l}{\textbf{\makecell{Engineering}}} & & \SetCell[r=1, c = 1]{l}{\cite{wiener1948cybernetics} \\ \cite{panteli2017metrics}} & \SetCell[r=1, c = 1]{l}{\cite{albert2000error} \\ \cite{gao2016universal}} & \SetCell[r=1, c = 1]{l}{\cite{bruneau2003framework} \\ \cite{francis2014metric} \\ \cite{panteli2017metrics}} & \SetCell[r=1, c = 1]{l}{\cite{rane2024artificial}} & \SetCell[r=1, c = 1]{l}{\cite{paries2017resilience}} & \SetCell[r=1, c = 1]{bg = gray!50} & \SetCell[r=1, c = 1]{l}{\cite{holling1996engineering} \\ \cite{francis2014metric} \\ \cite{ivanov2022viable}} \\
\hline
\SetCell[r=1,c=2]{l}{\textbf{\makecell{Ecology \& Social Systems}}} & & \SetCell[r=1, c = 1]{l}{\cite{holling1973resilience} \\  \cite{hashimoto1982reliability} \\ \cite{ives1995measuring}} & \SetCell[r=1, c = 1]{l}{\cite{bodin2009role}} & \SetCell[r=1, c = 1]{l}{\cite{ives1995measuring} \\ \cite{sundstrom2014transdisciplinary}} & \SetCell[r=1, c = 1]{l}{\cite{maillart2024computational}} & \SetCell[r=1, c = 1]{l}{\cite{berkes1998linking} \\ \cite{lebel2006governance} \\ \cite{mahajan2022participatory}} & \SetCell[r=1, c = 1]{l}{\cite{maillart2024computational}} & \SetCell[r=1, c = 1]{l}{\cite{folke2010resilience} \\ \cite{leslie2013response}\\ \cite{francis2014metric} \\ \cite{allen2018quantifying}}\\
\hline
\SetCell[r=1,c=2]{l}{\textbf{\makecell{Climate Change \& Disaster Management}}} & & \SetCell[r=1, c = 1]{l}{\cite{dawson2011agentbased} \\ \cite{panteli2017metrics} \\ \cite{aydin2018integration}} & \SetCell[r=1, c = 1]{l}{\cite{jacobs2017applying} \\ \cite{aydin2018integration}} & \SetCell[r=1, c = 1]{l}{\cite{bruneau2003framework}} & \SetCell[r=1, c = 1]{l}{\cite{arnold2004informationsharing} \\ \cite{aerts2018integrating}} & \SetCell[r=1, c = 1]{l}{\cite{dawson2011agentbased} \\ \cite{mahajan2022participatory}} & \SetCell[r=1, c = 1]{l}{\cite{moreno2018womens}} & \SetCell[r=1, c = 1]{l}{\cite{alexander2013resilience} \\ \cite{meerow2016defining} \\ \cite{adger2005socialecological} \\ \cite{leichenko2024climate}}\\
\hline
\SetCell[r=1,c=2]{l}{\textbf{\makecell{Psychology \& Neuroscience}}} & & \SetCell[r=1, c = 1]{l}{\cite{manikis2019computational} \\ \cite{ungar2021modeling}} & \SetCell[r=1, c = 1]{bg = gray!50} & \SetCell[r=1, c = 1]{l}{\cite{mcewen2004protection} \\ \cite{rakesh2019resilience}} & \SetCell[r=1, c = 1]{l}{\cite{klein2008performing} \\ \cite{bada2019cyber} \\\cite{kalisch2024neurobiology}} & \SetCell[r=1, c = 1]{l}{\cite{bada2019cyber}} & \SetCell[r=1, c = 1]{l}{\cite{rutter1987psychosocial} \\ \cite{masten2010disasters}} & \SetCell[r=1, c = 1]{l}{\cite{wagnild1993development} \\ \cite{luthar2000construct} \\ \cite{campbell2007psychometric} \\ \cite{feder2019biology}}\\
\hline
\SetCell[r=1,c=2]{l}{\textbf{\makecell{Public Health}}} & & \SetCell[r=1, c = 1]{l}{\cite{brauer2008compartmental} \\ \cite{luke2012systems}}  & \SetCell[r=1, c = 1]{l}{\cite{luke2012systems}} & \SetCell[r=1, c = 1]{l}{\cite{cutter2010disaster} \\ \cite{kruk2017building}} & \SetCell[r=1, c = 1]{l}{\cite{arnold2004informationsharing}} & \SetCell[r=1, c = 1]{l}{\cite{barasa2018what} \\ \cite{haldane2021health}} & \SetCell[r=1, c = 1]{l}{\cite{gibbs2013beyond} \\ \cite{aase2020resilience}} & \SetCell[r=1, c = 1]{l}{\cite{norris2007community} \\ \cite{barasa2018what}}\\
\hline
\SetCell[r=2, c=1]{l}{\vspace{-16.25em}\par \textbf{Management}} & \SetCell[r=1, c=1]{l}{\textbf{Organizational}} & \SetCell[r=1, c = 1]{bg = red!20} & \SetCell[r=1, c = 1]{bg = red!20} & \SetCell[r=1, c = 1]{l}{\cite{kuypers2018designing}} & \SetCell[r=1, c = 1]{l}{\cite{klein2008performing} \\ \cite{bada2019cyber} \\ \cite{almeida2020challenges}} & \SetCell[r=1, c = 1]{l}{\cite{lanzara1999transient} \\ \cite{westrum2004typology} \\ \cite{weick2011managing} \\ \cite{henfridsson2013generative} \\ \cite{barasa2018what} \\ \cite{bada2019cyber} \\\cite{boh2023building}} & \SetCell[r=1, c = 1]{bg = gray!50} &  \SetCell[r=1, c = 1]{l}{ \cite{westrum2004typology}\\ \cite{teece2007explicating} \\ \cite{vogus2007organizational} \\ \cite{bhamra2011resilience} \\ \cite{burnard2011organisational} \\ \cite{lengnick2011developing} \\ \cite{boin2013resilient} \\ \cite{henfridsson2013generative} \\ \cite{skopik2016problem} \\ \cite{dupont2019cyberresilience} \\ \cite{barasa2018what} \\ \cite{duchek2019organizational} \\ \cite{calder2021iso} \\\cite{hepfer2022heterogeneity} \\ \cite{he2022building}}\\
\cline[dashed]{2-9}
&  \SetCell[r=1, c=1]{l}{\textbf{Risk \& Cyber}} & \SetCell[r=1, c = 1]{bg = red!20} & \SetCell[r=1, c = 1]{bg = orange!20}{\cite{gillard2023efficient}} & \SetCell[r=1, c = 1]{l}{\cite{linkov2013resilience} \\ \cite{panteli2017metrics}} & \SetCell[r=1, c = 1]{l}{\cite{aerts2018integrating} \\ \cite{almeida2020challenges}} & \SetCell[r=1, c = 1]{l}{\cite{lostri2022shared} \\ \cite{boh2023building}} & \SetCell[r=1, c = 1]{l}{\cite{kuypers2018designing}} & \SetCell[r=1, c = 1]{l}{\cite{baselcommittee2011basel} \\ \cite{bjorck2015cyber} \\ \cite{woods2015four} \\ \cite{skopik2016problem} \\ \cite{linkov2019fundamental} \\ \cite{dupont2019cyberresilience} \\  \cite{hausken2020cyber}  \\ \cite{calder2021iso} \\ \cite{lostri2022shared}} \\
\end{longtblr}
\end{landscape}
\end{center}
\restoregeometry



Third, {\bf limited empirical validation and longitudinal grounding:} Field data from cybersecurity and risk practitioners indicates that resilience is typically assessed through compliance artifacts -- such as audits and certifications -- rather than by systematically analyzing adaptive behavior during disruption. Advanced tools such as predictive simulations or anomaly detection are piloted in innovation units but often remain decoupled from regulatory evaluation structures \autocite{skopik2016problem}. Informal practices -- parallel communication channels, crisis improvisation, or decentralized escalation -- remain critical to adaptive performance but largely invisible in existing assessments.\\

Similar dynamics are evident in healthcare and public health, where resilience is observed in practice yet only partially institutionalized in measurement systems \autocite{kruk2017building}. Moreover, longitudinal designs that track how resilience evolves across repeated disruptions remain rare. Few frameworks explicitly address the uncertainty, reversibility, and trade-offs inherent in resilience assessment \autocite{allen2018quantifying}. Disciplinary fragmentation further limits cross-fertilization, as methods from control theory or complex systems are seldom transferred into organizational research, and widely used practices such as awareness campaigns often fail to foster measurable behavioral adaptation \autocite{bada2019cyber}.\\
\\
Together, these gaps indicate that resilience research possesses substantial theoretical and methodological resources, yet lacks sufficient integration to support cumulative, organization-ready application. Without standardized, validated, and scalable metrics embedded in governance structures, resilience risks remaining a normative aspiration rather than a systematically cultivated capability. Mathematical approaches to dynamical systems—through basin stability, return dynamics, and perturbation response—provide promising foundations that could be adapted to trace organizational adaptation over time \autocite{krakovska2024resilience}. This is not solely a technical concern. For decision-makers, limited metric integration complicates investment prioritization, performance comparison, and accountability to boards and regulators. For scholars, it constrains intervention testing, replication, and theory-to-practice translation. Advancing resilience research therefore requires not additional definitions, but clearer integration of conceptual foundations, methodological tools, and empirical validation strategies. The next section addresses these gaps directly by clarifying operational mechanisms, advancing simulation-based measurement approaches, and outlining empirical pathways for longitudinal validation.

\subsection{Research in Organizational Resilience in Digital Environments}

Organizational resilience has traditionally been conceptualized as the capacity of firms to absorb shocks, adapt to disruption, and sustain core functions over time \autocite{duchek2020organizational,burnard2011organisational}. However, contemporary organizations operate through digitally mediated processes, interconnected infrastructures, and platform-based ecosystems \autocite{hepfer2022heterogeneity, janssen2006network}. As core activities become embedded in information systems, adaptive capacity is increasingly shaped by technological architectures, data flows, and algorithmic coordination mechanisms \autocite{hepfer2022heterogeneity,liu2022network}.\\

We therefore treat \emph{digital organizational resilience} as the contemporary sociotechnical specification of organizational resilience in digitally mediated environments (Duchek, 2020; Hepfer et al., [year]). Rather than constituting a distinct theoretical domain, digital organizational resilience clarifies how organizational resilience is enacted when adaptive capacity is distributed across digital infrastructures, institutional arrangements, and inter-organizational ecosystems (Janssen et al., 2006; Liu et al., 2022). In this view, resilience remains a dynamic capability (Duchek, 2020), yet its realization increasingly depends on digitally embedded coordination mechanisms and infrastructural interdependencies (Hepfer et al., [year]; Boin \& van Eeten, 2013).

Cyber resilience represents a particularly salient instantiation of this specification (Linkov et al., 2013). Adversarial digital threats expose the extent to which organizational resilience depends on information systems, cross-boundary coordination, and infrastructural interdependence (Dupont, 2019; Ibne Hossain et al., 2020). Studying resilience in cyber contexts thus renders visible the mechanisms through which organizational resilience becomes digitally structured and operationalized within interconnected sociotechnical systems (Linkov et al., 2013; Liu et al., 2022).

